%
% $Id: $
%
%
% Compilar a .pdf con LaTeX (pdflatex)
% Es necesario instalar Beamer (paquete latex-beamer en Debian)
%

%
% Gr�ficos:
% Los gr�ficos pueden suministrarse en PNG, JPG, TIF, PDF, MPS
% Los EPS deben convertirse a PDF (usar epstopdf)
%

\documentclass{beamer}
\usetheme{Warsaw}
\usebackgroundtemplate{\includegraphics[width=\paperwidth]{format/libresoft-bg-soft.png}}
\usepackage[spanish]{babel}
\usepackage[latin1]{inputenc}
\usepackage{graphics}
\usepackage{amssymb} % Simbolos matematicos

%\definecolor{libresoftgreen}{RGB}{162,190,43}
%\definecolor{libresoftblue}{RGB}{0,98,143}

%\setbeamercolor{titlelike}{bg=libresoftgreen}

%% Metadatos del PDF.
\hypersetup{
  pdftitle={Brief introduction to git},
  pdfauthor={Jesus M. Gonzalez-Barahona},
  pdfcreator={GSyC/LibreSoft, Universidad Rey Juan Carlos},
  pdfproducer=PDFLaTeX,
  pdfsubject={Case studies I},
}
%%

\begin{document}

\title{Brief introduction to git}
\subtitle{Case studies I / Master on libre software (URJC) \\
\url{http://master.libresoft.es}}
\author{Jesus M. Gonzalez-Barahona, Carlos Garc�a Campos}
\institute{jgb@gsyc.es \\
GSyC/LibreSoft, Universidad Rey Juan Carlos}

\date{September 2010}

\frame{
\maketitle
\begin{center}
\includegraphics[width=6cm]{format/gsyc-urjc}
\end{center}
}


% Si el titulo o el autor se quieren acortar para los pies de p�gina
% se pueden redefinir aqu�:
%\title{Titulo corto}
%\author{Autores abreviado}


%% LICENCIA DE REDISTRIBUCION DE LAS TRANSPAS
\frame{
~
\vspace{3cm}

\begin{flushright}
\copyright 2009-2010 Jesus M. Gonzalez-Barahona, Carlos Garc�a Campos. \\

Some rights reserved. \\
This document is distributed under the \\
Creative Commons Attribution-ShareAlike 3.0 licence, \\
available in \\
\url{http://creativecommons.org/licenses/by-sa/3.0}

The original version of this document is available at \\
\url{http://master.libresoft.es}
\end{flushright}
}
%%


\begin{frame}
\frametitle{Why are we introducing git?}

\begin{itemize}
\item One of the most used distributed version control systems
\item Originally written by Linus Torvalds for supporting Linux development
\item Now used by man projects
\item Available in most free software distributions
\item Available for many operating systems
\item Forges based on it available: Github, Gitorious, etc.
\item We will use it in the master's program
\end{itemize}

\begin{flushright}
\url{http://git-scm.com/} \\
\url{http://github.com/} \\
\url{http://gitorious.org/}
\end{flushright}

\end{frame}

\begin{frame}[fragile]
\frametitle{Starting to use git}

\begin{itemize}
\item Creating/initializing a new repository:
\begin{verbatim}
$ mkdir project
$ cd project
$ git init
\end{verbatim}
\item Only once for each repository
\item Everything (diffs for each commit) will be stored in it
\item No real difference between ``repository'' and ``working copy''
\end{itemize}
\end{frame}

\begin{frame}[fragile]
\frametitle{Normal workflow}

Committing:

\begin{verbatim}
$ git add file1 file2 ...
$ git diff --cached
$ git status
$ git commit 
\end{verbatim}

\begin{itemize}
\item add marks files to commit (adds them to the index)
\item diff to know what is about to be commited
\item status to get a general idea of the current state
\item commit really adds changes to files in the index
\item Special shortcut for committing
\begin{verbatim}
$ git commit -a
\end{verbatim}
\end{itemize}

The index is a snapshot of the working tree that represents the
content that will be affected by a commit
\end{frame}


\section{Exploring the history}

\begin{frame}[fragile]
  \frametitle{Exploring the history}

  \begin{itemize}
  \item {\it git log}
    {\scriptsize
\begin{verbatim}
commit ce6ddeac2200225b173c52e2509ffcc01b2b4a1d
Author: Carlos Garcia Campos <carlosgc@gnome.org>
Date:   Mon Nov 19 16:14:01 2007 +0100

    Initial commit
\end{verbatim}
    }


    \begin{itemize}
    \item Commit: commit id (sha1)
    \item Author: commit author (not the committer)
    \item Date: commit date
    \item Commit message: (a single line briefly 
      describing the change and, optionally, a blank line followed by
      one or more lines with a detailed description)
    \end{itemize}
  \end{itemize}

\end{frame}


\section{Operations between repositories}

\begin{frame}[fragile]
  \frametitle{Pulling}

  \begin{itemize}
  \item Cloning a repository
\begin{verbatim}
$ git clone <uri>
\end{verbatim}

  \item Getting updates
\begin{verbatim}
$ git fetch <uri>
$ git merge origin/master
\end{verbatim}

  \item Useful shortcut (commonly used)
\begin{verbatim}
$ git pull
\end{verbatim}

\item Updating a remote repository
\begin{verbatim}
$ git push
\end{verbatim}

  \end{itemize}
\end{frame}


\begin{frame}
\frametitle{Further references (git in detail)}

In this presentation, only a very brief tour to the basics...

To probe further:

\begin{itemize}
\item Class in the Development course in this master's
\item The Git Community Book \\
  \url{http://book.git-scm.com/}
\end{itemize}

\end{frame}

\end{document}
